\begin{beginnerstatement}[Kurz erklärt: Tauri und die Desktop-Schicht]

Tauri verwendet das vorhandene Web-Frontend und zeigt es in einer Webview,
also einem eingebetteten Web-Anzeigebereich, als Desktop-Anwendung an. Ein
Betriebssystem wie Windows, macOS oder Linux steuert den Computer; Rust bildet
hier die native Anwendungsschicht, die direkt damit arbeiten und native
Funktionen bereitstellen kann. Eine native Anwendung läuft
installiert auf dem Computer und nicht ausschließlich in einem normalen
Browser. Eine Capability beschreibt zusammen mit Berechtigungen, welche
nativen Funktionen erlaubt sind; die Content Security Policy (CSP) begrenzt
zusätzlich die nutzbaren Inhalte und Verbindungen. Ein Sidecar ist ein separates
Hilfsprogramm, das zusammen mit der Anwendung ausgeliefert werden kann, in der
Baseline aber nicht für FastAPI vorhanden ist. Eine Remote API läuft dagegen
auf einem entfernten Server und wird über das Netzwerk erreicht. Packaging
bündelt die Anwendung zu installierbaren Paketen, während Signing deren
Herkunft und Unverändertheit kryptografisch bestätigt. Ob Remote API, Sidecar
oder lokale Kommandos sowie Packaging und Signing verwendet werden, bleibt
hier eine Produktentscheidung.

\end{beginnerstatement}
